\documentclass[UTF8]{ctexart}
\usepackage{xeCJK}
\usepackage{geometry}
\geometry{left=2cm,right=2cm,top=2.5cm,bottom=2.5cm}
\setlength{\headheight}{12.7pt}

\usepackage{titlesec}
\titleformat{\section}
  {\normalfont\large\bfseries}
  {\thesection}
  {1em}
  {}
\titlespacing*{\section}
  {0pt}
  {3.5ex plus 1ex minus .2ex}
  {2.3ex plus .2ex}

\usepackage{amsmath}
\usepackage{amssymb}
\usepackage{amsthm}
\usepackage{enumitem}
\setlist[enumerate]{itemsep=0pt,topsep=0pt,parsep=0pt,leftmargin=0pt,itemindent=2.5em}
\setlist[itemize]{itemsep=0pt,topsep=0pt,parsep=0pt,leftmargin=0pt,itemindent=2.5em}
\usepackage{fancyhdr}
\pagestyle{fancy}
\fancyhf{}
\usepackage{graphicx}
\usepackage{hyperref}
\usepackage{indentfirst}
\usepackage{lmodern}


\begin{document}
\setlength{\abovedisplayskip}{1pt}
\setlength{\belowdisplayskip}{1pt}

\section{自然数的指数}

\hypertarget{sec:定义1.1}{}
\noindent\textbf{定义1.1 (自然数的指数)}

给定\href{../01 自然数/01 自然数#nameddest=sec:定义1.3}{自然数} $m$,
递归定义映射$\mathrm{exp}_m:\omega\to\omega$如下:
\begin{enumerate}
    \item $\mathrm{exp}_m(0)\overset{\mathrm{def}}{=}1$,
    \item $\mathrm{exp}_m(n^+)\overset{\mathrm{def}}{=}\mathrm{exp}_m(n)\cdot m$.
\end{enumerate}

接下来, 对任意的自然数$m,n$, 定义
\[
    m^n\overset{\mathrm{def}}{=}\mathrm{exp}_m(n),\\[-3pt]
\]
称为$m$的$n$\textbf{次方}.

\vspace{10pt}

显然任给非零自然数$n$有$0^n=0$, 任给自然数$n$有$1^n=1$.

\vspace{10pt}

\hypertarget{sec:定理1.1}{}
\noindent\textbf{定理1.1}

任给自然数$m,n,k$,
\begin{enumerate}
    \item $m^{n+k}=m^nm^k$,
    \item $m^{nk}=(m^n)^k$,
    \item $(mn)^k=m^kn^k$.
\end{enumerate}

\noindent{\kaishu 证明.}

\begin{enumerate}
    \item 任给自然数$m,n$, 对$k$归纳证明$m^{n+k}=m^nm^k$.
    
    \hspace{2.17em}
    首先$m^{n+0}=m^n=m^n1=m^nm^0$;
    
    \hspace{2.17em}
    其次任取自然数$k$, 假设$m^{n+k}=m^nm^k$, 那么
    \[
        m^{n+k^+}=m^{(n+k)^+}=m^{n+k}m=m^nm^km=m^nm^{k^+}.
    \]

    \item 任给自然数$m,n$, 对$k$归纳证明$m^{nk}=(m^n)^k$.
    
    \hspace{2.17em}
    首先$m^{n0}=m^0=1=(m^n)^0$;
    
    \hspace{2.17em}
    其次任取自然数$k$, 假设$m^{nk}=(m^n)^k$, 那么
    \[
        m^{nk^+}=m^{nk+n}=m^{nk}m^n=(m^n)^km^n=(m^n)^{k^+}.
    \]

    \item 任给自然数$m,n$, 对$k$归纳证明$(mn)^k=m^kn^k$.
    
    \hspace{2.17em}
    首先$(mn)^0=1=1\cdot 1=m^0n^0$;
    
    \hspace{2.17em}
    其次任取自然数$k$, 假设$(mn)^k=m^kn^k$, 那么
    \begin{equation*}
        (mn)^{k^+}=(mn)^kmn=m^kn^kmn=m^kmn^kn=m^{k^+}n^{k^+}.
    \tag*{\qedsymbol}\end{equation*}
\end{enumerate}





\newpage

然后列出一些指数与序的关系.

\vspace{10pt}

\hypertarget{sec:定理1.2}{}
\noindent\textbf{定理1.2}

任给自然数$m,n,k$,
\begin{enumerate}
    \item $n<2^n$,
    \item 如果$m<n$且$k>1$, 那么$k^m<k^n$,
    \item 如果$m\leqslant n$且$k\neq 0$, 那么$k^m\leqslant k^n$,
    \item 如果$m<n$且$k\neq 0$, 那么$m^k<n^k$,
    \item 如果$m\leqslant n$, 那么$m^k\leqslant n^k$.
\end{enumerate}

\noindent{\kaishu 证明.}

\begin{enumerate}
    \item 对$n$归纳证明$n<2^n$.
    
    \hspace{2.17em}
    首先$0<1=2^0$;
    其次任取自然数$n$, 假设$n<2^n$, 易证$1\leqslant 2^n$, 从而
    \[
        n^+=n+1<2^n+1\leqslant 2^n+2^n=2^n2=2^{n^+}.
    \]

    \item 记$q=n-m$, 显然$q\neq 0$. 易证$k^m\neq 0$且$1<k^q$, 所以
    \[        
        k^m<k^mk^q=k^{m+q}=k^n.
    \]

    \item 如果$k=1$, 那么$k^m=1\leqslant 1=k^n$;
    
    \hspace{2.17em}
    如果$k\neq 1$, 那么由第二条即得$k^m\leqslant k^n$.
    
    \item 假设$m<n$, 对$k$归纳证明$k\neq 0\to m^k<n^k$.
    
    \hspace{2.17em}
    首先$0=0$, 所以有$0\neq 0\to m^0<n^0$;
    
    \hspace{2.17em}
    其次任取自然数$k$, 假设$k\neq 0\to m^k<n^k$, 那么
    \[\left\{\begin{aligned}
        &k=0\implies&&m^{k^+}=m<n=n^{k^+},\\[-3pt]
        &k\neq 0\implies&&m^k<n^k\implies m^{k^+}=m^km<n^kn=n^{k^+}.
    \end{aligned}\right.\]

    \item 如果$k=0$, 那么$m^k=1\leqslant 1=n^k$.
    
    \hspace{2.17em}
    如果$k\neq 0$, 那么由第四条即得$m^k\leqslant n^k$. \hfill $\qedsymbol$
\end{enumerate}

% 这一堆可谓是毫无规律可言, 而且写了两个``易证''. 评价为拉完了.

\end{document}